% ─────────────────────────────────────────────────
%  CHAPITRE 2 – ANALYSE FONCTIONNELLE ET NON FONCTIONNELLE DES BESOINS
% ─────────────────────────────────────────────────
\chapter{Analyse fonctionnelle et non fonctionnelle des besoins}
\markboth{Analyse fonctionnelle et non fonctionnelle des besoins}{Analyse fonctionnelle et non fonctionnelle des besoins}

\section*{Introduction}
\addcontentsline{toc}{section}{Introduction}


Cette analyse précise les utilisateurs concernés, les fonctions attendues et
les exigences de qualité du système.

Dans de nombreux établissements, les informations liées aux clients et aux
investissements sont enregistrées dans plusieurs outils. C'est notamment le cas
des données KYC, des comptes bancaires, des portefeuilles et des ordres de
bourse. Les utilisateurs doivent donc consulter plusieurs applications. I-SIB
doit réunir ces informations sur une même plateforme.

Le chapitre présente d'abord le rôle du système et ses acteurs, puis les
fonctionnalités retenues. La création d'un prospect et le passage d'ordre sont
ensuite détaillés. Enfin, les exigences non fonctionnelles sont présentées.

% ─────────────────────────────────────────────────
\section{Rôle du système et acteurs}

Le rôle principal du système futur est d'offrir une plateforme bancaire intelligente qui
centralise la connaissance client et intègre les services SGI dans un même environnement.
Il doit organiser l'information de façon à assurer une continuité entre les parcours
bancaires et les parcours d'investissement, en répondant aux attentes spécifiques de chaque
type d'utilisateur.

Trois acteurs principaux ont été identifiés dans le périmètre de ce chapitre. Leurs
objectifs sont distincts mais complémentaires, et leur coexistence dans un même système
justifie la nécessité d'une architecture unifiée.

Le \textbf{client} est l'utilisateur final des services bancaires et d'investissement. Il
est au centre du dispositif : le système doit lui permettre d'accéder à ses informations
personnelles, à ses comptes, à ses produits bancaires et à son portefeuille
d'investissement. Il doit également pouvoir suivre la performance de ses actifs, consulter
les cotations et soumettre des demandes d'achat ou de vente de titres sans sortir de
l'écosystème bancaire.

Le \textbf{chargé de clientèle} est un utilisateur interne chargé du suivi commercial et
de l'accompagnement des clients.
Il attend du système qu'il facilite le suivi des clients et des prospects, tout en lui donnant accès aux informations dont il a besoin pour mener ses activités commerciales.
Il doit également permettre d'identifier plus facilement les clients nécessitant une attention particulière.

L'\textbf{équipe SGI} est un utilisateur interne spécialisé dans les opérations sur titres et la gestion de portefeuille. Elle utilise le système pour gérer les portefeuilles clients, traiter les ordres de bourse et suivre l'évolution des investissements à partir des informations de marché disponibles.

La plateforme regroupe les services de relation client et d'investissement. Les
droits attribués à chaque utilisateur déterminent les informations et les
fonctions auxquelles il peut accéder.

% ─────────────────────────────────────────────────
\section{Fonctionnalités du système}

Dans le périmètre étudié, I-SIB comprend les modules \textbf{Relation Client} et
\textbf{SGI} (Société de Gestion et d'Intermédiation). Le premier sert à gérer
les clients, les prospects et le suivi commercial. Le second permet de suivre
les portefeuilles et de traiter les ordres de bourse.


 La présente section décrit les fonctionnalités offertes par chacun de ces modules.

Pour chaque module, les principales fonctionnalités et les interactions avec les acteurs sont d'abord présentées. Une fonctionnalité jugée représentative est ensuite détaillée à travers une fiche descriptive puis illustrée afin de montrer les différentes étapes de son traitement.



% ─────────────────────────────────────────────────
\subsection{Module Relation Client}

Le module Relation Client regroupe l'ensemble des fonctionnalités permettant de centraliser
la connaissance client, de structurer le suivi commercial, de suivre les informations
KYC et de faciliter le traitement des demandes. La figure~\ref{fig:fonctionnalites_rc} présente la liste de ces fonctionnalités à travers une représentation mettant en évidence les interactions entre les acteurs et le système.


\begin{figure}[H]
    \centering
    \MemoireFigureTall{images/contenu/chapitre-2/use_case_relation_client.png}
    \caption{Liste des fonctionnalités du module Relation Client}
    \label{fig:fonctionnalites_rc}
\end{figure}

Dans la suite, nous allons détailler la fonctionnalité création d'un prospect.


% ─────────────────────────────────────────────────
\subsubsection{Création d'un prospect}

La création d'un prospect constitue le point d'entrée du pipeline commercial : le chargé
de clientèle enregistre un nouveau contact, que le système positionne au stade initial et
enrichit automatiquement d'un score de conversion calculé par intelligence artificielle.
La fiche~\ref{tab:rc_creation_prospect} en présente la description détaillée.



\begin{table}[H]
    \centering
    \caption{Fiche descriptive -- Création d'un prospect}
    \label{tab:rc_creation_prospect}
    \renewcommand{\arraystretch}{1.5}

    \begin{tabular}{|>{\bfseries}p{3.8cm}|p{10.5cm}|}
        \hline

        Acteurs &
        Chargé de clientèle \\ \hline

        Préconditions &
        \begin{minipage}[t]{\linewidth}\vspace{2pt}
            \begin{itemize}[leftmargin=*,topsep=0pt,itemsep=1pt,parsep=0pt]
                \item L'utilisateur est authentifié avec le rôle \textit{chargé de clientèle}
                \item Le module de gestion des prospects est accessible
            \end{itemize}
            \vspace{2pt}\end{minipage}
        \\ \hline

        Scénario nominal &
        \begin{minipage}[t]{\linewidth}\vspace{2pt}
            \begin{enumerate}[leftmargin=*,topsep=0pt,itemsep=2pt,parsep=0pt]
                \item Le chargé de clientèle renseigne les informations du prospect (nom, prénom, email, téléphone, source ; éventuellement société, valeur estimée et produits d'intérêt).

                \item Le système vérifie qu'aucun prospect existant n'est déjà associé à cette adresse email.

                \item Le prospect est créé au stade initial \textit{LEAD} et automatiquement rattaché au conseiller qui l'a enregistré. Un score initial lui est attribué.

                \item Le service de scoring calcule un score de conversion à partir des informations disponibles.

                \item Le score obtenu, la probabilité de conversion et le niveau de priorité sont ajoutés à la fiche du prospect.

                \item La fiche du prospect est communiquée avec les informations enregistrées.
            \end{enumerate}
            \vspace{2pt}\end{minipage}
        \\ \hline

        Scénario alternatif &
        \begin{minipage}[t]{\linewidth}\vspace{2pt}
            \textbf{A1 -- Adresse email déjà utilisée :} le système bloque la création et indique qu'un prospect utilise déjà cette adresse.
            \vspace{2pt}\end{minipage}
        \\ \hline

        Scénario d'exception &
        \textbf{E1 -- Service de scoring indisponible :}
        le prospect est quand même enregistré avec son score initial. L'incident est conservé dans le journal du système.
        \\ \hline

        Postconditions &
        Le prospect est enregistré, affecté au chargé de clientèle et placé au stade \textit{LEAD}. L'opération est inscrite dans le journal d'audit.
        \\ \hline

    \end{tabular}
\end{table}



La figure~\ref{fig:rc_prospects} représente le déroulement de cette fonctionnalité à
travers un diagramme de séquence.

\begin{figure}[H]
    \centering
    \MemoireFigureMedium{images/contenu/chapitre-2/rc_seq_prospects.png}
    \caption{Création d'un prospect et calcul automatique de son score de conversion}
    \label{fig:rc_prospects}
\end{figure}

% ─────────────────────────────────────────────────
\subsection{Module SGI}

Le module SGI intègre les services d'investissement au sein de la plateforme I-SIB. Il
permet au client de suivre son portefeuille et offre aux équipes SGI les outils nécessaires
pour superviser les actifs, traiter les ordres et analyser les performances. La
figure~\ref{fig:fonctionnalites_sgi} présente la liste des fonctionnalités de ce module.

\begin{figure}[H]
    \centering
    \MemoireFigureMedium{images/contenu/chapitre-2/use_case_sgi.png}
    \caption{Liste des fonctionnalités du module SGI}
    \label{fig:fonctionnalites_sgi}
\end{figure}

Dans la suite, nous allons détailler la fonctionnalité passage d'ordres.

% ─────────────────────────────────────────────────
\subsubsection{Passage d'ordres}

Le passage d'ordres constitue le lien opérationnel entre le client et l'équipe SGI. Il
permet au client d'exprimer une intention d'achat ou de vente, et à l'équipe SGI de
contrôler, valider ou rejeter la demande selon les règles métier en vigueur. La
fiche~\ref{tab:sgi_ordres} en présente la description détaillée.


\begin{table}[H]
\centering
\caption{Fiche descriptive -- Passage d'ordres}
\label{tab:sgi_ordres}
\renewcommand{\arraystretch}{1.5}
\begin{tabular}{|>{\bfseries}p{3.8cm}|p{10.5cm}|}
\hline
Acteurs & Client, Équipe SGI \\ \hline

    Préconditions &
    \begin{minipage}[t]{\linewidth}\vspace{2pt}
    \begin{itemize}[leftmargin=*,topsep=0pt,itemsep=1pt,parsep=0pt]
        \item Le client est authentifié et dispose d'un portefeuille
        \item Le titre concerné est référencé et sa cotation est disponible
        \item Le compte espèces associé au portefeuille est approvisionné
    \end{itemize}
    \vspace{2pt}\end{minipage} \\ \hline

    Scénario nominal &
    \begin{minipage}[t]{\linewidth}\vspace{2pt}
    \begin{enumerate}[leftmargin=*,topsep=0pt,itemsep=2pt,parsep=0pt]

        \item Le client saisit un ordre en précisant le titre, le type d'opération (achat ou vente) et la quantité souhaitée

        \item Le système vérifie que le portefeuille appartient au client et récupère les informations de cotation du titre

        \item L'ordre est enregistré avec le statut \textit{En attente}

        \item L'équipe SGI consulte les ordres en attente de traitement

        \item Après vérification, l'équipe SGI approuve l'ordre. Les contrôles nécessaires sont effectués et le statut devient \textit{Approuvé}

        \item L'équipe SGI exécute l'ordre. Les comptes concernés et le portefeuille sont mis à jour, les frais éventuels sont pris en compte et l'ordre passe au statut \textit{Exécuté}

        \item Le client consulte l'état mis à jour de son ordre et de son portefeuille

    \end{enumerate}
    \vspace{2pt}\end{minipage} \\ \hline

    Scénario alternatif &
    \begin{minipage}[t]{\linewidth}\vspace{2pt}
    \textbf{A1 -- Rejet de l'ordre :} l'équipe SGI refuse l'ordre en indiquant un motif. Le statut passe à \textit{Rejeté} et l'information est communiquée au client.
    \vspace{2pt}\end{minipage} \\ \hline

    Scénario d'exception &
    \begin{minipage}[t]{\linewidth}\vspace{2pt}

    \textbf{E1 -- Solde insuffisant :} lors de la validation, les fonds disponibles ne permettent pas d'exécuter l'opération. L'ordre n'est pas approuvé et reste au statut \textit{En attente}.

    \medskip

    \textbf{E2 -- Cotation indisponible :} les informations de marché nécessaires ne sont pas disponibles ; l'ordre ne peut pas être enregistré.

    \vspace{2pt}\end{minipage} \\ \hline

    Postconditions &
    L'ordre est enregistré avec son statut final. En cas d'exécution, le portefeuille ainsi que les comptes espèces et titres sont mis à jour. Les opérations réalisées sont conservées dans les journaux du système. \\ \hline

\end{tabular}
\end{table}



La figure~\ref{fig:sgi_ordres} représente le déroulement de cette fonctionnalité à
travers un diagramme de séquence.


\begin{figure}[H]
    \centering
    \MemoireFigureMedium{images/contenu/chapitre-2/sgi_seq_passage_ordres.png}
    \caption{Circuit de soumission et de traitement d'un ordre de bourse du client à l'exécution}
    \label{fig:sgi_ordres}
\end{figure}
% ─────────────────────────────────────────────────
\section{Exigences non fonctionnelles}

En plus des fonctionnalités, le système doit respecter plusieurs exigences non
fonctionnelles, qui traduisent les contraintes propres au domaine bancaire.
Côté \textbf{sécurité}, les fonctionnalités accessibles dépendent du rôle de
chacun. Un client ne voit que ses propres informations, tandis que les
collaborateurs disposent des droits liés à leurs missions. La
\textbf{confidentialité} va de pair, les données personnelles, bancaires, KYC et
d'investissement ne devant être consultées que par les personnes autorisées.
S'y ajoute la \textbf{traçabilité}, puisque chaque validation d'ordre,
modification de donnée ou changement de statut est enregistré, de façon à
pouvoir reconstituer l'historique des opérations et répondre aux contrôles. Le
système doit aussi rester \textbf{disponible}, pour que les utilisateurs
accèdent à tout moment aux informations dont ils ont besoin, et
\textbf{évolutif}, afin qu'un nouveau module ou une nouvelle fonctionnalité
puisse être ajouté sans tout reprendre. L'\textbf{ergonomie}, enfin, conditionne
l'adoption, car des écrans clairs et simples à utiliser limitent les erreurs de
manipulation.

L'ensemble de ces exigences devra être pris en compte lors de la conception et de la mise en œuvre de la solution.


% ─────────────────────────────────────────────────
\section*{Conclusion}
\addcontentsline{toc}{section}{Conclusion}

Ce chapitre a présenté les attentes auxquelles doit répondre I-SIB dans le périmètre des modules Relation Client et SGI. Il a précisé les utilisateurs concernés, le rôle attendu du système ainsi que les principales fonctionnalités à mettre en œuvre.

La création d'un prospect et le passage d'ordre ont été retenus pour illustrer plus concrètement le fonctionnement attendu. Le chapitre a également présenté les principales exigences non fonctionnelles auxquelles la solution devra répondre.

Le chapitre suivant est consacré à la conception de la solution proposée.
